\section{A.1. What the surgery maps can and cannot mean}
\subsection{The target in the manuscript}
On cote 67, scans 074--078, the normalization of the stable curve carries at every node a nonzero element of the tensor product of the two branch tangent lines. Scan 078 also uses its inverse as a perfect pairing. We denote the tensor by $q_e\in(T_{a_e}Y\otimes_{\C}T_{b_e}Y)^\times$. No bound on $q_e$ is imposed there. We first use this literal target. All curves in this section are complex, all marked points are distinct, and $2g-2+n>0$.

For a geodesic cut, the two induced boundary metrics agree, so the identification of the boundaries is an orientation-reversing isometry of circles. Canonical capping gives pointed embedded analytic discs. If instead arbitrary smooth curves are allowed without an isometric or toric gluing condition, the two canonical arclength actions need not agree; the asserted construction of $q_e$ then needs an additional definition. Our obstruction applies whenever the tensor is the one obtained from the actual pointed caps, independently of how the caps were selected.

\begin{lemma}[Known, with independent proof: the cap bound]\label{A1-cap-bound}
Let $f:\D\longrightarrow\C\setminus\{1\}$ be injective and holomorphic, with $f(0)=0$. Then $|f'(0)|\leq4$. Thus a pointed disc in $\mathbf P^1$ centred at $0$ and avoiding $1,\infty$ has conformal radius at most $4$ in the coordinate $z$.
\end{lemma}
\begin{proof}
We include the classical quarter-theorem argument. If $F(z)=z+a_2z^2+\cdots$ is univalent in $\D$, the exterior odd function
\[
G(w)=\frac{1}{\sqrt{F(w^{-2})}}=w-\frac{a_2}{2w}+\cdots\qquad(|w|>1)
\]
is univalent. The square root is chosen by its expansion at infinity; injectivity follows because equality of squares forces $w_1=\pm w_2$, while oddness excludes the minus sign. For an exterior univalent function $w+b_0+\sum_{k\geq1}b_kw^{-k}$, the area enclosed by the image of $|w|=r>1$ is
\[
\pi\left(r^2-\sum_{k\geq1}k|b_k|^2r^{-2k}\right)\geq0.
\]
The equality follows by integrating $\overline G\,dG/(2i)$, with uniformly convergent Laurent series on that circle. Letting $r\downarrow1$ gives the area theorem. Applied to $G$ it gives $|a_2|\leq2$. If $a\ne0$ is omitted by $F$, the function $F/(1-F/a)$ is again normalized and univalent, and has second coefficient $a_2+1/a$. Hence $|1/a|\leq4$. Apply this to $F=f/f'(0)$ and the omitted value $a=1/f'(0)$.
\end{proof}
\begin{proposition}[Known with derivation: cusp norms and decorations]\label{A2-cusp-norm}
For a compact pointed complex curve $(Y,p_1,\ldots,p_n)$ with $2g-2+n>0$, the complete hyperbolic metric on $Y\setminus\{p_i\}$ defines canonical Hermitian norms on $T_{p_i}Y$. If $q_i$ is its cusp coordinate, then
\[
 ds^2=\frac{|dq_i|^2}{|q_i|^2\log^2|q_i|},\qquad
 \|v\|_{p_i}=|dq_i(v)|.
\]
The coordinate is unique up to a constant of modulus one. A horocycle of circumference $h_i$ has coordinate radius $\exp(-2\pi/h_i)$ wherever embedded. Penner's unrestricted decorations, understood as horocycles in the cusp cover and allowing their immersed images, form a principal $(\R_{>0})^n$-bundle. The map which assigns to $X$ the capped curve and decorations of circumferences $\ell_i(X)$ identifies its positive-length bundle with the one in Theorem \ref{A4-grafting}.
\end{proposition}
\begin{proof}
Normalize a primitive parabolic stabilizer of a cusp in the upper half-plane to $\zeta\mapsto\zeta+1$. Then $q=\exp(2\pi i\zeta)$ gives the displayed metric. Any two such normalized isometries differ by $\zeta\mapsto\zeta+a$ with $a\in\R$, giving precisely the unit scalar ambiguity and hence the intrinsic norm. The horocycle $\operatorname{Im}\zeta=y$ has length $1/y$; its $q$-radius is $e^{-2\pi y}$. Circumferences give independent positive real coordinates on the decoration fibre. Now apply Theorem \ref{A4-grafting}. This is Wolpert, \emph{Cusps and the family hyperbolic metric}, Section 1, Definition 1, together with the definition of Penner's decorated space. It applies in genus zero as well; the absence of a canonical smooth metric on the whole sphere causes no obstruction to a canonical cusp tangent norm. The norm does not assert that a unit tangent disc is realized by an embedded cap.
\end{proof}

\begin{theorem}[Proved: metric comparison after removing normal collars]\label{A2-collar}
Let $(M,g)$ be a compact hyperbolic surface with geodesic boundary and negative Euler characteristic. Remove pairwise disjoint normal collars of widths $d_i>0$ from selected boundary components, without changing the topological type, and call the remaining surface $M'$. Let $g'$ be its conformally compatible hyperbolic metric with geodesic boundary. Then $g'=e^{2u}g$ on $M'$ with $u>0$. More precisely $u$ is the unique smooth solution of
\[
\Delta_g u=e^{2u}-1\quad\hbox{on }M',\qquad
\partial_{n}u=\begin{cases}\tanh d_i&\hbox{on a new boundary},\\0&\hbox{on a retained boundary},\end{cases}
\]
where $\Delta=\operatorname{div}\operatorname{grad}$ and $n$ is the outward unit normal of $M'$. The new boundary corresponding to a former boundary of length $L_i$ has length strictly greater than $L_i\cosh d_i$ in $g'$.
\end{theorem}
\begin{proof}
Uniformization supplies $g'$ and its smooth extension to the analytic boundary. The curvature transformation gives $\Delta_g u=e^{2u}-1$. In normal coordinates at an old boundary, $g=dr^2+\cosh^2(r)ds^2$, with $r$ increasing into $M$. At $r=d_i$ the outward normal of the retained surface is $-\partial_r$, and the curvature vector of its unit tangent is $-\tanh(d_i)\partial_r$. Its outward geodesic curvature is therefore $\tanh d_i$. Under conformal scaling, this signed curvature becomes $e^{-u}(k_g-\partial_nu)$, giving the boundary condition.

If $u$ had a negative minimum in the interior, the equation would give $\Delta_g u<0$ there, a contradiction. If a negative minimum occurred on the boundary, write the equation as $\Delta_g u-c(x)u=0$ with $c(x)=(e^{2u}-1)/u>0$ (continuously set $c=2$ at zero). The boundary point lemma gives $\partial_nu<0$ for a nonconstant solution at such a minimum, contradicting the nonnegative prescribed derivative. A negative constant cannot solve the equation. Thus $u\geq0$. The strong maximum principle and the same boundary lemma imply $u>0$ everywhere, since a positive collar width excludes the identically zero solution.

For uniqueness, if $u_1,u_2$ solve the same problem, their difference $v$ satisfies $\Delta_gv=c(x)v$ with $c(x)>0$ and homogeneous Neumann boundary data. Integration by parts gives $-\int|\nabla v|^2=\int c v^2$, so $v=0$. Finally the old metric length of $r=d_i$ is $L_i\cosh d_i$ and the strict conformal factor inequality increases it. Thus the question of 67/021 has an intrinsic boundary-value answer; the old metric is not simply its restriction after the collar is removed.
\end{proof}

\begin{proposition}[Refuted, with rechecked extremal-length certificate: the orange slices]\label{A4-orange-extremal}
Let $Q_0=\{|z|\leq1,\ \operatorname{Re}z\leq1/2\}$ in the sphere with centres $0,1,\infty$. The canonical cap $D_0$ for a generalized pant of boundary lengths $(32,0,0)$ is not contained in $Q_0$. In particular the literal model in the letter to Bers, which permits the other two caps to degenerate to points, is false.
\end{proposition}
\begin{proof}
We recheck the extremal-length argument credited in the input to Michel Hua's project. Suppose $D_0\subset Q_0$. For $z\in Q_0$,
\[
 |z+1/2|^2=|z|^2+\operatorname{Re}z+1/4\leq7/4.
\]
Since $\sqrt7/2<4/3<7/5<3/2$, the annulus
\[
 A=\{4/3<|z+1/2|<7/5\}
\]
lies in $P=\mathbf P^1\setminus(D_0^\circ\cup\{1,\infty\})$, and its core is freely homotopic in $P$ to the boundary of $D_0$. For this curve family $\gamma$, monotonicity under inclusion of the annulus gives
\[
\operatorname{Ext}_P(\gamma)\leq\operatorname{Ext}_A(\gamma)
=\frac{2\pi}{\log(21/20)}.
\]
The hyperbolic metric on $P$ has area $2\pi$ by Gauss--Bonnet, and the shortest curve in the boundary class has length $32$. In the defining supremum of extremal length it therefore gives
\[
\operatorname{Ext}_P(\gamma)\geq\frac{32^2}{2\pi}.
\]
These bounds contradict $1024\log(21/20)>4\pi^2$. For an exact elementary verification, $\log(21/20)=\int_{20}^{21}dt/t>1/21$, while $1024/21>4(22/7)^2>4\pi^2$. The annulus extremal length is obtained by the logarithmic coordinate, which gives a Euclidean rectangle with width $\log(21/20)$ and periodic height $2\pi$. Cauchy--Schwarz gives the matching upper bound for every conformal density. The two-cusp pant exists for every positive remaining boundary length. Every geometric and numerical hypothesis in the contradiction is thus satisfied.
\end{proof}

\begin{lemma}[Proved: the tensor produced by a toric seam]\label{A1-tensor}
Let $f:\overline\D\to D_a\subset Y$ and $h:\overline\D\to D_b\subset Y$ parametrize two pointed analytic discs, with centres $a,b$. An orientation-reversing toric identification of their boundaries has the form $f(\zeta)\sim h(c/\zeta)$, $|\zeta|=1$, where $|c|=1$. Its tangent tensor is
\[
q=c\,df_0(\partial_\zeta)\otimes dh_0(\partial_\eta).
\]
In any chosen branch coordinates its modulus is the product of the two conformal radii.
\end{lemma}
\begin{proof}
The annular continuation of the boundary identification is $\zeta\eta=c$. Applying the derivatives to the two tangent factors gives the displayed tensor. Replacing $f(\zeta)$ by $f(\alpha\zeta)$ and $h(\eta)$ by $h(\beta\eta)$ changes $c$ to $c/(\alpha\beta)$ and the derivative factors to $\alpha df_0,\beta dh_0$, so the tensor is independent of the two rotations. Taking absolute values proves the last assertion.
\end{proof}

\begin{theorem}[Refuted: the global isomorphisms in A.1 and L2]\label{A1-global-refutation}
With the unrestricted nonzero-tensor target of cote 67/075--078, the proposed maps cannot be isomorphisms. Already for $(g,n)=(0,4)$ the unbordered map is not surjective on complex objects. The same obstruction occurs for every stable type with $3g-3+n>0$, for the unbordered map and for the bordered map wherever it is defined by canonical capping. It is an obstruction to coarse surjectivity as well as to equivalence of stacks or groupoids.
\end{theorem}
\begin{proof}
For $(0,4)$ take two copies of $\mathbf P^1$, identify their points $0$ to form a node, and use $1,\infty$ on each as the four labelled markings. This is a connected stable curve of arithmetic genus zero with four markings. Put
\[
q=17\,\partial_z|_0\otimes\partial_w|_0.
\]
It is a valid nonzero nodal tangent tensor. Any preimage under the asserted surgery would give, on each normalized sphere, an embedded pointed cap about $0$ avoiding $1$ and $\infty$. Lemmas \ref{A1-cap-bound} and \ref{A1-tensor} would imply $|q|\leq4\cdot4=16$, contrary to its definition. The markings fix the chosen normalization, so this contradiction is invariant under isomorphism.

For every stable $(g,n)$ of positive complexity, a pants decomposition gives a connected stable trivalent graph with $n$ legs, first Betti number $g$, and $3g-3+n>0$ edges. Replace each vertex by a sphere with its three flags at $0,1,\infty$. At each flag use a local projective coordinate sending that flag to $0$ and the other two flags to $1,\infty$. The ambiguity of interchanging the latter two is $z\mapsto z/(z-1)$, whose derivative at zero has modulus one. Thus each flag has a well-defined tangent norm, preserved by all automorphisms of the graph curve. Choose a nonzero tensor of norm $17$ at every edge. A preimage would require all these norms to be at most $16$. This is impossible, including for a self-node, such as the one-node genus-one curve with one marking.

For the bordered target also choose admissible external discs. These exist: at each trivalent vertex take a generalized hyperbolic pant with cusp ends at the incident nodes and arbitrary positive lengths at the incident external legs; canonically cap just its actual boundaries. The resulting external discs are admissible on the sphere with its node branches left as marked punctures. The node tensors in the stated target are independent extra data, so choose their norms to be $17$ as above. Any hypothetical preimage would still have embedded node caps avoiding the other special points. This gives the same contradiction.
\end{proof}

\begin{proposition}[Proved: the zero-complexity edge case]\label{A1-zero}
For $(g,n)=(0,3)$ there is no essential stable cut. The unbordered groupoid with labelled markings is a point. The bordered capping functor is an equivalence with pointed spheres equipped with an admissible three-disc system, when admissible has its definition from 67/065--066 and arrows retain those discs.
\end{proposition}
\begin{proof}
The sphere with three ordered points has a unique projective normalization and no nontrivial automorphism fixing all three. For the bordered assertion, remove the three open discs. Admissibility says exactly that canonical arclength capping reconstructs the same pointed discs. Conversely, removing the caps from a canonically capped pant recovers the pant. Conformal isomorphisms restrict and extend uniquely across the analytic seams. These are inverse functors, not an identification of arbitrary radii with discs.
\end{proof}

\begin{proposition}[Proved: the corrected surgery equivalence]\label{A1-retained-discs}
Fix a stable graph. Enlarge the target by retaining admissible pointed caps at every branch of every node, as well as at the external markings, and a toric boundary identification at every node. Require the two hyperbolic boundary lengths at a node to be equal. Then geodesic cutting and canonical capping give an equivalence of groupoids with hyperbolic bordered surfaces cut along that graph. The nodal tensor alone is a forgetful invariant of this equivalence.
\end{proposition}
\begin{proof}
Remove all open caps. Each component has its unique hyperbolic metric with geodesic boundary. By admissibility its boundary circle actions are the given ones. A toric orientation-reversing identification between circles of equal circumference is an isometry. In signed normal distance $r$ the collar metric is $dr^2+\cosh^2(r)\,ds^2$; therefore the two collars glue to a smooth curvature-minus-one metric with the seam a geodesic. Its conformal structure gives the required welded surface. Cutting along the indicated seams recovers the original pieces and caps. Both operations carry isomorphisms to isomorphisms and are mutually inverse. If there are punctures, they retain their complete cusp neighborhoods throughout. The uniqueness of the conformal structure can also be read from analytic reflection across these seams.
\end{proof}

\subsection{What local plumbing does prove}
Hubbard--Koch, Theorem 9.11, establishes a homeomorphism onto an open neighborhood for sufficiently small plumbing parameters; the sentence immediately following the theorem explicitly says that the plumbing map is not globally injective. Their Theorem 10.1 gives the complex manifold structure and the universal marked family. Their Section 13 compares the quotient with the analytification of the Deligne--Mumford compactification. These statements do not provide an inverse to the unrestricted-tensor map refuted above.

\begin{proposition}[Known with derivation: the intrinsic first-order parameter]\label{A1-first-order}
For an ordinary node with branches $a,b$ on a smooth normalization $Y$, the one-dimensional space of local first-order smoothings is canonically $T_aY\otimes T_bY$. Given analytic coordinates $x,y$ at the branches, the family $xy=t$ represents this identification by $t\,\partial_x\otimes\partial_y$. Coordinate choices enter at higher order. For a fixed stable complex curve, sufficiently small such parameters together with deformation parameters for the normalized pointed components give local analytic plumbing charts, modulo its finite automorphism group.
\end{proposition}
\begin{proof}
For $A=\C[[x,y]]/(xy)$, the hypersurface presentation gives
\[
\operatorname{Ext}^1_A(\Omega_A,A)=A/(x,y)=\C.
\]
Explicitly, the relation $y\,dx+x\,dy=0$ yields the cokernel of $(u,v)\mapsto yu+xv$. Replacing $x,y$ by coordinates with linear terms $ax,by$ replaces $t$ by $ab\,t$, while the tensor frame changes by $(ab)^{-1}$. Thus the tensor is intrinsic. Nonlinear coordinate terms are irrelevant modulo the square of the deformation ideal, but can change actual finite plumbing parameters. The asserted local charts and finite quotient are exactly the application of Hubbard--Koch, Sections 8--10, to the normalization with its branch markings; stability makes the automorphism group finite. This is a local analytic assertion and supplies no global algebraic isomorphism.
\end{proof}

\section{A.2. Capping and its dependence on the other boundaries}
\begin{proposition}[Proved: projective connection of a capped pant]\label{A2-schwarzian}
Let a hyperbolic pant have boundary lengths $\ell_0,\ell_1,\ell_\infty>0$, and canonically cap it at the three boundaries. Normalize the resulting sphere by placing the cap centres at $0,1,\infty$. Set
\[
a_j=\frac12\left(1+\frac{\ell_j^2}{4\pi^2}\right).
\]
The hyperbolic developing map extends as a complex projective structure to the punctured capped sphere, whose Schwarzian in the coordinate $z$ is
\[
T(z)=\frac{a_0}{z^2}+\frac{a_1}{(z-1)^2}
       +\frac{a_\infty-a_0-a_1}{z(z-1)}.
\]
If $u$ is any centred unit-disc coordinate on the cap at $0$, write $u(z)=c z(1+bz+O(z^2))$. Then
\[
b=\frac{a_0+a_1-a_\infty}{2a_0}.
\]
\end{proposition}
\begin{proof}
Choose arclength $s$ on a boundary of length $\ell$, and a conformal collar coordinate $v=s+iy$ with boundary $y=0$. After an isometry of the upper half-plane, a hyperbolic developing map is $i\exp(v)$. The flat cylinder attached by arclength has punctured-disc coordinate $u=\exp(2\pi i v/\ell)$. Thus the continued projective developing map is $i u^{-i\ell/(2\pi)}$, with dilation monodromy. Its Schwarzian is $(1+\ell^2/(4\pi^2))/(2u^2)$. This gives a meromorphic projective connection with the stated double-pole coefficients at the three centres. On $\mathbf P^1$ a rational function with only these double poles and the required double-pole behavior at infinity is uniquely of the displayed form: a difference would be a holomorphic quadratic differential with at most simple poles at three points, whose degree is $-4+3<0$.

By the Schwarzian chain rule,
\[
T(z)=a_0\left(\frac{u'(z)}{u(z)}\right)^2+\{u,z\}.
\]
Here $u'/u=z^{-1}+b+O(z)$ and $\{u,z\}$ is holomorphic at zero. The simple-pole coefficient on the right is $2a_0b$, and that on the left is $a_0+a_1-a_\infty$. This proves the formula. Multiplication of $u$ by a rotation changes $c$ but not $b$.
\end{proof}

\begin{theorem}[Refuted: the independence asked in 67/068]\label{A2-independence}
In a canonically capped pant the cap at $0$ need not depend only on $\ell_0$. In particular, the pants with length triples $(2\pi,2\pi,2\pi)$ and $(2\pi,4\pi,2\pi)$ have different caps at $0$ in the common normalization $0,1,\infty$.
\end{theorem}
\begin{proof}
Hyperbolic pants exist and are unique for both positive length triples. Proposition \ref{A2-schwarzian} gives respectively $b=1/2$ and $b=5/4$. If the two pointed caps were the same domain, their centred disc coordinates would differ by a unit scalar, so their $b$ coefficients would agree. They do not. The corrected dependence is explicit in that proposition: at fixed $\ell_0$, the second jet already records $\ell_1^2-\ell_\infty^2$.
\end{proof}

\begin{corollary}[Refuted: fixed concentric orange-slice caps]\label{A4-orange-jets}
There are no three fixed pointed discs $Q_0,Q_1,Q_\infty$ whose independently concentric subdiscs describe all canonically capped pants. Thus the literal model in 67/019--020 is false, independently of its additional spherical description.
\end{corollary}
\begin{proof}
If $U:Q_0\to\D$ is a centred uniformization, the unit-disc coordinate on the concentric subdisc $U^{-1}(\lambda\D)$ is $U/\lambda$. Its coefficient $b$ in Proposition \ref{A2-schwarzian} is independent of $\lambda$. The two actual pants in Theorem \ref{A2-independence} have different $b$.
\end{proof}

\subsection{The canonical bundle furnished by grafting}
\begin{theorem}[Known with derivation: capping is infinite grafting]\label{A4-grafting}
Let $S$ be a marked compact oriented surface of type $(g,n)$ with $n\geq1$ and negative Euler characteristic. Let $\mathcal T_B^{\geq0}(S)$ classify complete hyperbolic structures with each labelled end either a geodesic boundary of length $\ell_i>0$ or a cusp of length $\ell_i=0$, and with no nodes in the interior. Let $\mathcal T_{g,n}$ be its punctured Teichm\"uller space. Canonical capping gives a mapping-class-equivariant homeomorphism
\[
\Phi=(\operatorname{cap},\ell):\mathcal T_B^{\geq0}(S)
\longrightarrow\mathcal T_{g,n}\times[0,\infty)^n.
\]
Its restriction to each boundary-type stratum is a real-analytic diffeomorphism. In particular the positive-length part is a trivial principal $(\R_{>0})^n$-bundle over $\mathcal T_{g,n}$ by boundary lengths, not by all possible tangent radii.
\end{theorem}
\begin{proof}
The conformal cylinder $(\R/\ell_i\Z)\times[0,\infty)$ is the punctured unit disc through $u=\exp(2\pi i(s+iy)/\ell_i)$. Adding its missing point gives exactly the pointed disc attached by normalized arclength in 67/065--066. Thus the cap operation is Mondello's $\operatorname{gr}_\infty$, not just an analogy with it. Apply Mondello, Section 5.4, Theorem 5.4, restricted to the finite, non-nodal strata. His length vector is the same $\ell$ and his marking convention is the one obtained by collapsing boundary circles to the punctures. At zero total length his blown-up decoration directions are forgotten; the continuous maps descend to the displayed product. For all positive lengths the theorem supplies both the inverse and real analyticity, and on each mixed zero/positive stratum it supplies the same assertion. Hence the comparison transfers exactly the claimed structure. A section is $Y\mapsto\Phi^{-1}(Y,(1,\ldots,1))$.
\end{proof}

\begin{corollary}[Proved: an intrinsic action and retraction]\label{A4-action}
The multiplicative monoid $[0,\infty)^n$ acts continuously and intrinsically on $\mathcal T_B^{\geq0}(S)$ by
\[
\lambda\cdot X=\Phi^{-1}\bigl(\operatorname{cap}(X),(\lambda_i\ell_i(X))_i\bigr).
\]
It commutes with the mapping class action (and with a simultaneous permutation of end labels). On the positive-length part its group of units acts freely and transitively on each capping fibre. The zero vector gives the canonical capping retraction. The maps $X\mapsto\Phi^{-1}(\operatorname{cap}(X),(1-t)\ell(X))$ give a strong deformation retraction onto the punctured locus.
\end{corollary}
\begin{proof}
All assertions follow by conjugating coordinatewise multiplication and straight contraction by the specified homeomorphism $\Phi$. Its definition uses only canonical capping and the intrinsic hyperbolic boundary lengths. It therefore makes no choice of a pants decomposition. It need not agree with the operation that holds unrelated Fenchel--Nielsen coordinates fixed while multiplying selected lengths.
\end{proof}
